\section{Conclusion and future work}\label{sec:concl}
\subsection*{Scope and limitations}
The present work formalises only the first layer of the APEIRON hierarchy.
Extending the categorical semantics to the full seventeen‑layer architecture is a long‑term programme that will require proving adjunctions, handling probabilistic co‑activation, and eventually incorporating continuous dynamics.
We view the current paper as the necessary algebraic substrate upon which such higher‑order structures can be built.
The fact that the decomposition operators are implemented using practical proxies (e.g., zlib for Kolmogorov complexity) does not diminish the categorical result: the free category $\mathcal{A}$ is defined relative to any given family of operators, and the separatedness theorem holds for any concrete instantiation.
Improving the quality of the operators will refine the category, but not invalidate it.

We have presented a rigorous categorical foundation for the multi‑axial atomicity of the APEIRON Framework.
Defining the Apeiron category as the free category over the decomposition graph and proving that atomicity is equivalent to being a separated object provides a clean, compositional interpretation of the framework's most fundamental layer.
The inter‑layer functors then extend this compositional structure to the entire architecture.

Future work includes:
\begin{itemize}
\item Developing the full categorical semantics for Layers 2–7, proving adjunctions that model emergent behaviour.
\item Extending the category‑theoretic description to the qualitative dimensions and relational density fields.
\item Formalising the autopoietic and world‑building layers using higher‑order categorical structures (e.g., monads, comonads).
\item Incorporating the Coq and Lean proofs of atomicity into a machine‑verified categorical proof assistant.
\end{itemize}

The present paper lays the categorical groundwork that makes all these steps feasible.

\section*{Acknowledgments}
The author thanks the anonymous reviewers for their constructive feedback.